\documentclass[10pt,a4paper]{amsart}
\usepackage[margin=19mm]{geometry}
\usepackage{amsmath,amssymb,mathtools}
\usepackage[scr=rsfs,cal=euler]{mathalfa}
\usepackage{xfrac}
\usepackage[unicode,hidelinks,bookmarks=false]{hyperref}
\newcommand{\ie}{i.e.\ }
\newcommand{\cf}{cf.\ }
\newcommand{\resp}{respectively\ }
\newcommand{\Subsection}{\subsection}
\providecommand{\nobreakdash}{\nobreak\hbox{-}}
\setlength{\emergencystretch}{2em}
\multlinegap0pt
\usepackage{bm}
\usepackage{bbm}
\usepackage{dsfont}
\usepackage{xspace}
\usepackage{cancel}
\usepackage{mathtools}
\usepackage{amsthm}
\makeatletter
\@ifundefined{claim}{\newtheorem{claim}{Claim}}{}
\@ifundefined{theorem}{\newtheorem{theorem}{Theorem}}{}
\makeatother
\title[Complexity and algorithms for proper conflict-free coloring ]{Complexity and algorithms for proper conflict-free coloring in graphs}
\author{Dinabandhu Pradhan, Vaishali Sharma}
\date{Source: arXiv:2608.10874v2 (CC BY 4.0)}
\begin{document}
\maketitle

\noindent\emph{Source and licence.} Complexity and algorithms for proper conflict-free coloring in graphs. Version arXiv:2608.10874v2 (CC BY 4.0), distributed under the Creative Commons Attribution 4.0 International licence, as recorded in the arXiv licence metadata for this exact version and shown on the abstract page https://arxiv.org/abs/2608.10874v2. Full rights notice: licenses/arxiv-2608.10874-CC-BY-4.0.txt.\par\smallskip

\noindent\textit{Editorial scope.} Counted unit: the Theorem whose source label is \texttt{split} in the article (source0.tex lines 917--923, source label \texttt{split}), delivered together with the article's own complete original proof of it (source0.tex lines 924--967, one \texttt{proof} environment of 44 lines). No mathematics is written, repaired or re-derived: statement and proof are the article's own text line for line. Required context retained uncounted: none. Every definition, display and label of those lines is retained unchanged. Omitted from this package and not redistributed: 1--916, 968--1081. Nothing inside a retained excerpt is omitted or altered: every retained body is delivered exactly as the article's lines are written, so no omission receipt of any kind accompanies this package. Citations: the retained excerpts carry no citation command at all, so no bibliography entry of the article is transcribed and no reference is redistributed. Reference adaptation: none was needed and none was made; every reference inside the delivered unit resolves inside it, so no unresolved reference or citation is produced. Printed slips kept literal: no printed word, symbol, hypothesis or number of the counted statement or of its proof was altered. Layout and class: the article's own class is \texttt{article} with packages including color, ulem, tikz, todonotes, epsfig, enumerate, float, amsmath, amsfonts, amsbsy, amssymb, amsthm, mathrsfs, lineno; the wrapper is the lane's pinned \texttt{amsart} editorial wrapper at 10pt on A4, which restates the theorem-like environments the retained text needs and adds the article's own macro definitions that the retained text needs (none), with extra wrapper packages (bm, bbm, dsfont, xspace, cancel, mathtools). The delivered numbering is the unit's own. Assets: no retained span contains a figure environment, a graphics-inclusion command, a PDF-page inclusion, a table or a TikZ picture; no image file is copied into this package, the package declares no asset dependency and no image file is redistributed with it. Licence: version arXiv:2608.10874v2 of this article is distributed under the Creative Commons Attribution 4.0 International licence, as recorded in the arXiv licence metadata for this exact version and shown on the abstract page https://arxiv.org/abs/2608.10874v2. A full rights notice is delivered at licenses/arxiv-2608.10874-CC-BY-4.0.txt. Characters: every retained excerpt is pure ASCII, so the delivered engine, XeLaTeX with its default Latin Modern text font, reports no missing character.
\begin{theorem}\label{split}
	Let $G$ be a split graph with $V(G) = C \cup I$, where $C$ is a maximal clique with $|C|\geq 2$ and $I$ is an independent set. Then 
	\begin{enumerate}
		\item $\chi_{pcf}(G) = \omega(G)$ if and only if $G$ is special.
		\item $\chi_{pcf}(G) = \omega(G)+1$ if and only if $G$ is not special.
	\end{enumerate}
\end{theorem}

\begin{proof}
	
	To prove (a) and (b), we first prove the following claim.
	\begin{claim}\label{bound_pseudo-split}
		If $G$ is not special, then $\chi_{pcf}(G)\geq \omega(G)+1$.
	\end{claim}
	\begin{proof}
		For the sake of contradiction, assume that there exists a PCF $\omega(G)$-coloring, say $f$.
		Let $G$ be a split graph which is not special. Let $C = \{c_1, c_2, \ldots, c_k\}$. 
		First, assume that $v_i$ does not exist. 
		Consequently, there exists a vertex $x_i\in I$ such that $N(x_i) = C\setminus\{c_i\}$ for each vertex $c_i\in C$. Since $C$ is a clique, each vertex of $C$ must be assigned distinct colors under $f$. Without loss of generality, assume that $f(c_i) = i$ for each $c_i\in C$. Clearly, $x_i$ must be assigned color $i$ for $f$ to be a proper $k$-coloring. Therefore, $f(x_i) = i$ for every vertex $x_i$, $i\in [k]$.
		It is easy to see that none of the vertices of $C$ has a unique color in its neighborhood. This is a contradiction to the fact that there exists a PCF $\omega(G)$-coloring. Thus, $\chi_{pcf}(G)\geq \omega(G)+1$.
		
		\medskip
		
		Now, assume that $v_i$ exists but $v_j$ does not exist.  Since $C$ is a clique, each vertex of $C$ must be assigned distinct colors under $f$. Without loss of generality, assume that $f(c_i) = i$ for each $c_i\in C$. 
		Since $v_j$ does not exist, there exists a vertex $x_j\in N(v_i)\cap I$ such that $N(x) = C\setminus\{v_j\}$ for each vertex $c_j\in C\setminus\{v_i\}$. Note that $x_j$ must be assigned the color $j$ for $f$ to be a proper $k$-coloring. 
		Therefore, each color appears twice in $N(v_i)$. Thus, $v_i$ does not have a unique color in its neighborhood. 
		This is a contradiction to the fact that there exists a PCF $\omega(G)$-coloring. Thus, $\chi_{pcf}(G)\geq \omega(G)+1$.
		This completes the proof of Claim~\ref{bound_pseudo-split}.
	\end{proof}
	
	Now, we prove statements (a) and (b).
	
	\medskip
	
	(a) The necessary part follows by Claim~\ref{bound_pseudo-split}. To prove the sufficiency, let $G$ be a special split graph. Then, either $|C| = 1$ or there exist vertices $v_i, v_j \in C$ such that 
	\begin{itemize}
		\item if $x\in I$ is not adjacent to $v_i$, then $x$ has a non-neighbor in $C\setminus\{v_i\}$.
		\item if $y\in N(v_i)\cap I$ is not adjacent to $v_j$, then $y$ has a non-neighbor in $C\setminus\{v_j\}$.
	\end{itemize}
	
	Let $C = \{c_1, c_2, \ldots, c_k\}$.
	Let $f$ be a coloring such that $f(c_i) = i$ for each $c_i \in C$. Since $C$ is a maximal clique, every vertex of the set $I$ is non-adjacent to some vertex of $C$. Now color every vertex $v\in I\setminus N(v_i)$ with a color from the set $\{1, 2, \ldots, k\}$ different than the colors assigned to its neighbors and $f(v_i)$. Since there is no vertex $x\in I$ such that $f(v_i)$ is the only color available from the set $[k]$ for $x$ due to properness, we can color each vertex $x\in I\setminus N(v_i)$ with a color distinct from its neighbors and $f(v_i)$. 
	Finally, color every vertex $v\in N(v_i)\cap I$ with a color from the set $\{1, 2, \ldots, k\}$ different than the colors assigned to its neighbors and $f(v_j)$. Since there is no vertex $x'\in N(v_i)\cap I$ such that $f(v_j)$ is the only color available from the set $[k]$ for $x'$ due to properness, we can color each vertex $x\in N(v_i)\cap I$ with a color distinct from its neighbors and $f(v_j)$. 
	Clearly, $f$ uses $\omega(G)$ colors. Since every vertex is assigned a color different from the colors assigned to its neighbors, $f$ is a proper coloring. Since $N(v)$ is a clique for each vertex $v\in I$, each vertex $v\in I$ has a unique color in $N(v)$. Since the color $f(v_i)$ does not appear on any vertex of the set $I$, $f(v_i)$ is a unique color in $N(x)$ for each vertex $x\in C\setminus\{v_i\}$. Since the color $f(v_j)$ does not appear on any vertex $v'\in N(v_i)\cap I$, $f(v_j)$ is a unique color of $v_i$ in $N(v_i)$. Since we need at least $\omega(G)$ colors for proper coloring of $G$ and $f$ uses $\omega(G)$ colors, $f$ is an optimal PCF $\omega(G)$-coloring of the graph $G$. Thus, $\chi_{pcf}(G) = \omega(G)$.
	
	\medskip
	
	(b) We provide a PCF coloring with $\omega(G)+1$ colors. Let $f$ be a coloring of $G$ such that $f(c_i) = i$ for every vertex $c_i\in C$ and $f(x) = k+1$ for every vertex $x\in I$. Since adjacent vertices are assigned different colors, $f$ is a proper coloring. Note that $C$ has at least two vertices.
	Clearly, every vertex $c_i\in C\setminus\{c_1\}$ has a unique color $f(c_1)$ in $N(c_i)$. The vertex $c_1$ has a unique color $f(c_2)$ in $N(c_1)$.
	Since $N(v)$ is a clique for each $v\in I$, each vertex $v\in I$ has a unique color in $N(v)$. Therefore, $f$ is a PCF coloring with $\omega(G)+1$ colors. Consequently, $\chi_{pcf}(G)\leq\omega(G)+1$. By Claim~\ref{bound_pseudo-split}, $\chi_{pcf}(G)\geq\omega(G)+1$. 
	Thus, $\chi_{pcf}(G) = \omega(G)+1$.
\end{proof}

\end{document}
